\clearpage
\begin{proofbox}
    \textbf{\textcolor{CProof}{要证明的到底是什么？}}
    先证明所有函数值\textbf{不会超过}两个端点的范围；
    再证明这两个端点之间的\textbf{每一个数都能取到}。
    只做第一步不够，因为“不会超过”不等于“每个都能取到”。

    \medskip
    \textbf{\textcolor{CProof}{方法一：从大小关系和图像理解（门槛最低）}}
\begin{center}\begin{tikzpicture}[font=\small,>=stealth,line join=round]
\draw[gray!45,rounded corners=3pt] (0,0) rectangle (14.2,9.4);
\node at (7.1,8.8) {连续单调曲线的全部高度，由两端点夹住};

\begin{scope}[shift={(2.3,2.2)},x=2cm,y=1.3cm]
\fill[blue!7] (0,1) rectangle(2,3.8);
\draw[->] (-.3,0)--(2.6,0) node[right]{$x$};\draw[->] (0,-.2)--(0,4.5) node[above]{$y$};
\draw[teal!80!black,very thick,domain=0:2,samples=70] plot(\x,{1+\x+.1*\x^3});
\draw[dashed] (1.3,0)--(1.3,2.5197)--(0,2.5197);\draw[dashed] (2,0)--(2,3.8)--(0,3.8);
\foreach \x/\y in {0/1,1.3/2.5197,2/3.8}{\fill[blue] (\x,\y) circle(1.7pt);}
\node[below] at (0,0){$a$};\node[below] at (1.3,0){$x$};\node[below] at (2,0){$b$};\node[left] at (0,1){$f(a)$};\node[left] at (0,2.5197){$f(x)$};\node[left] at (0,3.8){$f(b)$};
\end{scope}
\node[align=left,anchor=west] at (8.4,5.1){$a\le x\le b$\\[10pt] $f(a)\le f(x)\le f(b)$\\[12pt] 单调：确定上下界\\[8pt] 连续：保证中间不缺值\\[8pt] 闭区间：两端可取到};
\node[text width=13.2cm,align=center] at (7.1,.65) {示意递增情形；递减时，最小值在右端、最大值在左端。};
\end{tikzpicture}
\end{center}

    假设 $f(x)$ 在 $[a,b]$ 上连续且单调递增。
    \begin{description}[style=nextline,leftmargin=2.8em,labelsep=0.5em,itemsep=0.48em]
        \item[\textbf{第一步：先排出 $x$ 的顺序。}]
            取区间内任意一个 $x$，都有 $a\le x\le b$。
        \item[\textbf{第二步：让函数值跟着排列。}]
            因为函数递增，所以
            \[
                f(a)\le f(x)\le f(b).
            \]
            这表明函数值不可能小于 $f(a)$，也不可能大于 $f(b)$。
        \item[\textbf{第三步：看有没有漏掉的高度。}]
            图像从 $(a,f(a))$ 连续走到 $(b,f(b))$，中途没有跳跃。
            因此两端高度之间每个高度都能走到。
        \item[\textbf{第四步：写出值域。}]
            两端点本身又在定义域里，所以 $f(a)$ 和 $f(b)$ 也能取到。
            故值域为 $[f(a),f(b)]$。
    \end{description}
    如果函数单调递减，只需把第二步的不等号方向调过来，得到
    $f(b)\le f(x)\le f(a)$，其值域为 $[f(b),f(a)]$。

    \medskip
    \end{proofbox}
\clearpage
\begin{proofbox}
\textbf{\textcolor{CProof}{方法二：挑出任意一个中间值来证明（高中常规）}}
\begin{center}\begin{tikzpicture}[font=\small,>=stealth,line join=round]
\draw[gray!45,rounded corners=3pt] (0,0) rectangle (14.2,9.4);
\node at (7.1,8.8) {减去任意目标高度，问题变成寻找横轴交点};

\begin{scope}[shift={(2.9,4.1)},x=2cm,y=1.3cm]
\draw[->] (-.4,0)--(2.6,0) node[right]{$x$};\draw[->] (0,-1.5)--(0,2.4) node[above]{$g(x)$};
\draw[teal!80!black,thick,domain=0:2,samples=70] plot(\x,{\x+.1*\x^3-1.1});
\fill[blue] (0,-1.1) circle(2pt);\fill[blue] (2,1.7) circle(2pt);\fill[orange!90!black] (1,0) circle(2pt);\draw[dashed] (2,0)--(2,1.7);
\node[below left] at (0,-1.1){$g(a)<0$};\node[above right] at (2,1.7){$g(b)>0$};\node[below right] at (1,0){$c$};\node[below] at (2,0){$b$};\node[above left] at (0,0){$a$};
\end{scope}
\node[align=left,anchor=west] at (8.7,5.1){$f(a)<t<f(b)$\\[10pt] $g(x)=f(x)-t$ 连续\\[10pt] $g(a)<0<g(b)$\\[10pt] 存在 $c$：$g(c)=0$\\[8pt] 即 $f(c)=t$};
\node[text width=13.2cm,align=center] at (7.1,.65) {从横轴下方连续走到上方，必经过横轴；这个交点实现了目标值。};
\end{tikzpicture}
\end{center}

    仍设 $f(x)$ 在 $[a,b]$ 上连续且单调递增。
    \begin{description}[style=nextline,leftmargin=2.8em,labelsep=0.5em,itemsep=0.48em]
        \item[\textbf{第一步：先确定不会超出端点。}]
            由单调性，仍有 $f(a)\le f(x)\le f(b)$。
        \item[\textbf{第二步：随意挑一个目标高度。}]
            任取满足 $f(a)<t<f(b)$ 的数 $t$。
            我们要证明：一定存在某个 $x$，使得 $f(x)=t$。
        \item[\textbf{第三步：把问题变成寻找零点。}]
            令 $g(x)=f(x)-t$，则
            \[
                g(a)=f(a)-t<0,\qquad g(b)=f(b)-t>0.
            \]
            由于 $f$ 连续，减去一个固定的数 $t$ 后，$g$ 也连续。
            图像从横轴下方走到横轴上方，中间必然经过横轴，
            所以存在 $c\in(a,b)$，使得 $g(c)=0$。
        \item[\textbf{第四步：还原成函数值。}]
            $g(c)=0$ 就是 $f(c)=t$。
            由于 $t$ 是随意挑选的，说明两端之间的每个数都能取到；
            再加上两个端点，值域便是 $[f(a),f(b)]$。
    \end{description}
    递减时可以把 $g(a)$ 和 $g(b)$ 的正负关系对调，证明完全相同。
    如果 $f(a)=f(b)$ 且整个区间单调，则函数始终等于这个数，
    值域为只含一个数的区间 $[f(a),f(a)]$。

    \medskip
    \end{proofbox}
\clearpage
\begin{proofbox}
\textbf{\textcolor{CProof}{方法三：从答案倒着找 $x$（适合可以直接解出 $x$ 的函数）}}
\begin{center}\begin{tikzpicture}[font=\small,>=stealth,line join=round]
\draw[gray!45,rounded corners=3pt] (0,0) rectangle (14.2,9.4);
\node at (7.1,8.8) {任取区间内的目标值，都能反解出合法自变量};

\begin{scope}[shift={(2.5,2.1)},x=2cm,y=1.25cm]
\draw[->] (-.2,0)--(2.8,0) node[right]{$x$};\draw[->] (0,-.2)--(0,4.8) node[above]{$y$};\draw[blue,thick,domain=0:2] plot(\x,{\x*\x});
\draw[orange!90!black,dashed] (0,2.25)--(1.5,2.25)--(1.5,0);\fill[orange!90!black] (1.5,2.25) circle(2pt);\fill[blue] (2,4) circle(2pt);\fill[blue] (0,0) circle(2pt);
\node[left] at (0,2.25){$y$};\node[below] at (1.5,0){$\sqrt y$};\node[below] at (2,0){$2$};\node[left] at (0,4){$4$};\draw[dashed] (2,0)--(2,4)--(0,4);
\end{scope}
\node[align=left,anchor=west] at (8.5,5.1){$f(x)=x^2,\ x\in[0,2]$\\[12pt] 任取 $y\in[0,4]$\\[10pt] $x=\sqrt y\in[0,2]$\\[10pt] $f(\sqrt y)=y$};
\node[text width=13.2cm,align=center] at (7.1,.65) {这里的水平线与竖直线给出“从输出找输入”的实际对应。};
\end{tikzpicture}
\end{center}

    这种办法不用先依赖图像：\textbf{给出一个目标函数值，反过来构造能得到它的 $x$。}
    以 $f(x)=x^2$、$0\le x\le2$ 为例。
    \begin{description}[style=nextline,leftmargin=2.8em,labelsep=0.5em,itemsep=0.48em]
        \item[\textbf{第一步：确定函数值不会超出范围。}]
            因为 $0\le x\le2$，所以 $0\le x^2\le4$。
        \item[\textbf{第二步：任取范围内的数。}]
            随便给定 $0\le y\le4$，我们想找 $x$ 让 $x^2=y$。
        \item[\textbf{第三步：直接把 $x$ 找出来。}]
            取 $x=\sqrt{y}$。由于 $0\le y\le4$，可知 $0\le\sqrt y\le2$，
            这个 $x$ 确实在题目规定的区间内。
        \item[\textbf{第四步：代回核对。}]
            $f(\sqrt y)=(\sqrt y)^2=y$。
            所以 $[0,4]$ 内每个 $y$ 都有对应的 $x$，值域恰为 $[0,4]$。
    \end{description}
    \textbf{说明：}第三种是在常见的具体函数上直接证明“中间没有漏值”，
    尤其适合一次函数和能通过开平方解出 $x$ 的二次函数；
    它不是对任意函数都能直接套用的通用方法。

    \medskip
    \textbf{\textcolor{CProof}{最后再检查一个反面情况}}
    若只知道单调，不知道连续，前两种方法中“中间每个高度都能经过”这一步就没有保证。
    因此原来仅凭“单调 $+$ 闭区间”就认定完整值域是端点区间的证明，并不充分。
\end{proofbox}
